\section{Manejo de series temporales}

\subsection{Conversión de fechas (\texttt{to\_datetime})}

Cuando se importan \textit{Data sets}, estos suelen traer las fechas como texto. Con \texttt{pd.to\_datetime} se transforman en objeto tipo \texttt{datetime64}. Listos para operaciones temporales.

\begin{pycode}{Conversion de texto a fecha}
import pandas as pd

datos = {
    "Fecha": ["2024-01-01", "2024-01-05", "2024-02-10", "2024-02-15", "2024-03-01"],
    "Ventas": [100, 120, 90, 150, 200]
}

df = pd.DataFrame(datos)

# Lo convertimos en fecha
df["Fecha"] = pd.to_datetime(df["Fecha"])
print(df.dtypes) # Retorna:
# Fecha     datetime64[us]
# Ventas             int64
# dtype: object
\end{pycode}



\subsection{Indexación por fecha}

Una fecha puede ser establecida como índice con \texttt{set\_index('columna de fechas')}, para luego poder indexar desde esta.

\begin{pycode}{Indexación por índice}
import pandas as pd

datos = {
    "Fecha": ["2024-01-01", "2024-01-05", "2024-02-10", "2024-02-15", "2024-03-01"],
    "Ventas": [100, 120, 90, 150, 200]
}

df = pd.DataFrame(datos)
df["Fecha"] = pd.to_datetime(df["Fecha"])

# Fechas como índice
df = df.set_index("Fecha")
print(df) # Retorna:
#             Ventas
# Fecha             
# 2024-01-01     100
# 2024-01-05     120
# 2024-02-10      90
# 2024-02-15     150
# 2024-03-01     200

print(df.loc["2024-01-5"]) # Retorna:
# Ventas    120

print(df.loc["2024-01-05":"2024-03-01"]) # Retorna:
#             Ventas
# Fecha             
# 2024-01-05     120
# 2024-02-10      90
# 2024-02-15     150
# 2024-03-01     200    
\end{pycode}

\begin{nota}
    Expresiones como: \texttt{2024} devuelve el año entero y \texttt{2024-01} devuelve el mes de ese año.
\end{nota}


\subsection{Resampleo (\texttt{.resample()})}

La función \texttt{.resample()} re-agrupa datos por periodo (día, mes, año) y le aplica una función (mean(), count()\dots).

Para reagrupar datos se utiliza como frecuencia:
\begin{itemize}
    \item \texttt{'D'}: día
    \item \texttt{'W'}: Semana
    \item \texttt{'ME'}: Mes
    \item \texttt{'QE'}: Trimestre
    \item \texttt{'YE'}: Año
\end{itemize}

\begin{pycode}{Resampleo por periodo}
import pandas as pd

datos = {
    "Fecha": ["2024-01-01", "2024-01-05", "2024-02-10", "2024-02-15", "2024-03-01"],
    "Ventas": [100, 120, 90, 150, 200]
}

df = pd.DataFrame(datos)
df["Fecha"] = pd.to_datetime(df["Fecha"])

df = df.set_index("Fecha")

print(df.resample('QE').count()) # Retorna:
#             Ventas
# Fecha             
# 2024-03-31       5

print(df.resample('ME').count()) # Retorna:
#             Ventas
# Fecha             
# 2024-01-31       2
# 2024-02-29       2
# 2024-03-31       1
\end{pycode}






\section{Optimización de Recursos}

\textit{Pandas} es muy bueno, pero su flexibilidad tiene un costo en memoria RAM. Por defecto, la librería es \textit{generosa} y asigna el espacio máximo posible a cada dato para evitar errores, pero esto puede volver pesados sus procesos.

\subsection{Eficiencia de memoria (\textit{Downcasting})}

Cuando usted carga un archivo, \textit{Pandas} suele asignar tipos de datos estándar como \texttt{int64} o \texttt{float64}. Esto significa que cada número reserva 64 bits de memoria, incluso si el número es un simple ``5''.

\subsubsection{Reducción de tipos numéricos}

Podemos utilizar \texttt{.astype()} para encoger nuestros datos a variantes que ocupen menos bits, siempre y cuando el valor quepa en ese nuevo envase:

\begin{table}[h]
\centering
\begin{tabular}{|l|l|l|}
\hline
\textit{Tipo} & \textit{Rango de valores} & \textit{Uso recomendado} \\ \hline
\texttt{int8} & -128 a 127 & IDs cortos, cantidades pequeñas. \\ \hline
\texttt{int16} & -32,768 a 32,767 & Stock de depósitos, días del año. \\ \hline
\texttt{int32} & Más de 9 ceros & Conteos estándar. \\ \hline
\texttt{float32} & Decimales con 7 dígitos & Precios y tasas de interés. \\ \hline
\end{tabular}
\caption{Variantes de tipos de datos para optimización}
\end{table}

\begin{pycode}{Downcasting manual}
import pandas as pd

# Supongamos que tiene una columna de 'Edad' o 'Mes'
# No hace falta usar 64 bits para un número que no pasa de 100

df['Mes'] = df['Mes'].astype('int8')

# Si tiene precios con pocos decimales, float32 es más que suficiente
df['Precio'] = df['Precio'].astype('float32')
\end{pycode}

\subsubsection{El tipo \texttt{category}: El ahorro definitivo}

Como mencionamos anteriormente, si una columna de texto se repite constantemente (como ``País de Destino'' o ``Precio''), convertirla a \texttt{category} puede reducir el uso de memoria en un 90\%. En lugar de guardar la palabra ``Argentina'' mil veces, \textit{Pandas} la guarda una sola vez y le asigna un número pequeño a las repeticiones.

\begin{pycode}{Optimización por categorías}
Verificamos el uso de memoria antes y después

print(df['Pais'].memory_usage(deep=True))

df['Pais'] = df['Pais'].astype('category')

print(df['Pais'].memory_usage(deep=True)) # Notará una reducción masiva
\end{pycode}

\begin{nota}
Puede usar \texttt{df.info(memory\_usage='deep')} para ver cuántos Megabytes está consumiendo su tabla en tiempo real.
\end{nota}


Utilice funciones \textit{vectorizadas} siempre que pueda. Solo utilice bucles o las funciones de \texttt{.apply()} y \texttt{.map()} para lógicas personalizadas.




\subsection{Estrategias para Datasets de Gran Volumen}

Cuando el volumen de datos excede la memoria RAM disponible, es necesario abandonar la carga tradicional y adoptar un enfoque de procesamiento incremental y almacenamiento eficiente.

\subsection{Lectura incremental (\textit{chunksize})}

El parámetro \texttt{chunksize} convierte la carga en un iterador. Es decir le indica a \textit{Pandas} que no cargue todo el archivo en la memoria, sino que lo convierta en un objeto iterable que entrega una cantidad fija de filas por cada vuelta.

\begin{pycode}{Procesamiento por trozos}
import pandas as pd

# Procesamos el archivo en bloques de 100,000 filas
for chunk in pd.read_csv("datos_grandes.csv", chunksize=100000):
    # Realizamos una agregación parcial para no saturar la memoria
    print(f"Promedio parcial de salario: {chunk['Salario'].mean()}")
\end{pycode}

\subsection{Optimización desde la carga (\textit{dtype} y \textit{usecols})}

Puede definir los tipos de datos y las columnas necesarias durante la lectura. (Antes que entren a la memoria)

\begin{itemize}
\item \texttt{usecols}: Carga solo las columnas estrictamente necesarias para el análisis. Menos columnas significan menos uso de RAM.
\item \texttt{dtype}: Asigna tipos de datos reducidos (\texttt{int8}, \texttt{float32}, \texttt{category}) desde el inicio.
\end{itemize}

\begin{pycode}{Carga quirúrgica}
import pandas as pd

# Cargamos solo 3 columnas con tipos de datos optimizados
df = pd.read_csv("datos.csv",
    usecols=["ID", "Edad", "Ciudad"],
    dtype={"Edad": "int8", "Ciudad": "category"})
\end{pycode}

\subsection{Formatos de almacenamiento de alto rendimiento}

El formato \texttt{.csv} es universal pero ineficiente para archivos grandes (es lento de leer y no preserva los tipos de datos). Para análisis profesional, se utilizan formatos columnares como Parquet o Feather.

\begin{table}[h]
\centering
\begin{tabular}{|l|l|l|}
\hline
\textit{Formato} & \textit{Ventaja Principal} & \textit{Uso Ideal} \\ \hline
\texttt{.csv} & Universalidad. & Intercambio de datos pequeños. \\ \hline
\texttt{.parquet} & Compresión. & Almacenamiento a largo plazo. \\ \hline
\texttt{.feather} & Velocidad récord. & Almacenamiento temporal. \\ \hline
\end{tabular}
\caption{Comparativa de formatos de archivo en Pandas}
\end{table}

\begin{pycode}{Uso de Parquet}
import pandas as pd
import pyarrow

#Exportar a formato optimizado (requiere la librería pyarrow)
df.to_parquet("datos.parquet", index=False)

# Lectura 
df = pd.read_parquet("datos.parquet")
\end{pycode}

\subsection{Delegación de carga (SQL y Motores Externos)}

Si sus datos residen en una base de datos SQL. Antes de traer los datos, utilice el motor de la base de datos para realizar filtros (\texttt{WHERE}) y agregaciones (\texttt{GROUP BY}) antes de importar los resultados a \textit{Pandas}.